\chapter{Node- und Item-Callbacks}
\label{ch:items/callbacks}

Luanti setzt beim Modding stark auf Callbacks.\index{callback@Callback}
Ein Callback ist eine Funktion, die du einer API übergibst und die
aufgerufen wird, wenn ein Ereignis eintritt. Du kannst zum Beispiel in
einer Node-Definition eine Funktion
\texttt{on\_\allowbreak{}punch}\index{callback@Callback!on_punch@\texttt{on\_\allowbreak{}punch}}
angeben, die aufgerufen wird, wenn ein Spieler den Node schlägt. Es gibt
auch globale Callbacks wie
\texttt{core.\allowbreak{}register\_\allowbreak{}on\_\allowbreak{}punchnode}\index{core@\texttt{core}!register_on_punchnode@\texttt{core.\allowbreak{}register\_\allowbreak{}on\_\allowbreak{}punchnode}},
die Ereignisse für alle Nodes empfangen.

\section{Item-Callbacks}\label{item-callbacks}

Hat ein Spieler einen Node, ein Craftitem\index{craftitem@Craftitem}
oder ein Werkzeug im Inventar,\index{inventar@Inventar} kann er damit
bestimmte Ereignisse auslösen:

\begin{longtable}[]{@{}
  >{\raggedright\arraybackslash}p{(\linewidth - 4\tabcolsep) * \real{0.3333}}
  >{\raggedright\arraybackslash}p{(\linewidth - 4\tabcolsep) * \real{0.2821}}
  >{\raggedright\arraybackslash}p{(\linewidth - 4\tabcolsep) * \real{0.3846}}@{}}
\toprule\noalign{}
\begin{minipage}[b]{\linewidth}\raggedright
Callback
\end{minipage} & \begin{minipage}[b]{\linewidth}\raggedright
Standardbelegung
\end{minipage} & \begin{minipage}[b]{\linewidth}\raggedright
Standardwert
\end{minipage} \\
\midrule\noalign{}
\endhead
\bottomrule\noalign{}
\endlastfoot
\texttt{on\_\allowbreak{}use} & Linksklick & \texttt{nil} \\
\texttt{on\_\allowbreak{}place} & Rechtsklick auf einen Node &
\texttt{core.\allowbreak{}item\_\allowbreak{}place} \\
\texttt{on\_\allowbreak{}secondary\_\allowbreak{}use} & Rechtsklick ins
Leere &
\texttt{core.\allowbreak{}item\_\allowbreak{}secondary\_\allowbreak{}use}
(tut nichts) \\
\texttt{on\_\allowbreak{}drop} & Q &
\texttt{core.\allowbreak{}item\_\allowbreak{}drop} \\
\texttt{after\_\allowbreak{}use} & Abbauen eines Nodes & \texttt{nil} \\
\end{longtable}

\subsection{on\_use}\label{on_use}

Hat ein Item einen
\texttt{on\_\allowbreak{}use}\index{callback@Callback!on_use@\texttt{on\_\allowbreak{}use}}-Callback,
kann es nicht mehr zum Abbauen von Nodes verwendet werden. Ein typischer
Einsatz von \texttt{on\_\allowbreak{}use} sind Lebensmittel:

\begin{Shaded}
\begin{Highlighting}[]
\VariableTok{core}\OperatorTok{.}\NormalTok{register\_craftitem}\OperatorTok{(}\StringTok{"mymod:mudpie"}\OperatorTok{,} \OperatorTok{\{}
    \VariableTok{description} \OperatorTok{=} \StringTok{"Alien{-}Matschkuchen"}\OperatorTok{,}
    \VariableTok{inventory\_image} \OperatorTok{=} \StringTok{"myfood\_mudpie.png"}\OperatorTok{,}
    \VariableTok{on\_use} \OperatorTok{=} \VariableTok{core}\OperatorTok{.}\NormalTok{item\_eat}\OperatorTok{(}\DecValTok{20}\OperatorTok{),}
\OperatorTok{\})}
\end{Highlighting}
\end{Shaded}

Die Zahl, die du
\texttt{core.\allowbreak{}item\_\allowbreak{}eat}\index{core@\texttt{core}!item_eat@\texttt{core.\allowbreak{}item\_\allowbreak{}eat}}
übergibst, ist die Anzahl der
Lebenspunkte,\index{lebenspunkte@Lebenspunkte} die der Verzehr heilt.
Jedes Herz des Spielers ist zwei Lebenspunkte wert. Ein Spieler hat
normalerweise bis zu 10 Herzen, also 20 Lebenspunkte.

\texttt{core.\allowbreak{}item\_\allowbreak{}eat()} ist eine Funktion,
die eine Funktion zurückgibt; diese wird als
\texttt{on\_\allowbreak{}use}-Callback eingetragen. Der Code oben ist
also gleichbedeutend mit:

\begin{Shaded}
\begin{Highlighting}[]
\VariableTok{core}\OperatorTok{.}\NormalTok{register\_craftitem}\OperatorTok{(}\StringTok{"mymod:mudpie"}\OperatorTok{,} \OperatorTok{\{}
    \VariableTok{description} \OperatorTok{=} \StringTok{"Alien{-}Matschkuchen"}\OperatorTok{,}
    \VariableTok{inventory\_image} \OperatorTok{=} \StringTok{"myfood\_mudpie.png"}\OperatorTok{,}
    \VariableTok{on\_use} \OperatorTok{=} \KeywordTok{function}\OperatorTok{(...)}
        \ControlFlowTok{return} \VariableTok{core}\OperatorTok{.}\NormalTok{do\_item\_eat}\OperatorTok{(}\DecValTok{20}\OperatorTok{,} \KeywordTok{nil}\OperatorTok{,} \OperatorTok{...)}
    \KeywordTok{end}\OperatorTok{,}
\OperatorTok{\})}
\end{Highlighting}
\end{Shaded}

Wenn du verstehst, dass \texttt{item\_\allowbreak{}eat} einfach eine
Funktion zurückgibt, kannst du sie abwandeln und komplexeres Verhalten
einbauen, etwa einen eigenen Sound abspielen.

\subsection{on\_place und
on\_secondary\_use}\label{on_place-und-on_secondary_use}

\texttt{on\_\allowbreak{}place}\index{callback@Callback!on_place@\texttt{on\_\allowbreak{}place}}
wird aufgerufen, wenn der Spieler auf einen Node zeigt,
\texttt{on\_\allowbreak{}secondary\_\allowbreak{}use}\index{callback@Callback!on_secondary_use@\texttt{on\_\allowbreak{}secondary\_\allowbreak{}use}},
wenn er das nicht tut.

Beide Callbacks werden für alle Item-Arten aufgerufen.
\texttt{on\_\allowbreak{}place} ist standardmäßig
\texttt{core.\allowbreak{}item\_\allowbreak{}place}\index{core@\texttt{core}!item_place@\texttt{core.\allowbreak{}item\_\allowbreak{}place}};
diese Funktion ruft entweder den
\texttt{on\_\allowbreak{}rightclick}\index{callback@Callback!on_rightclick@\texttt{on\_\allowbreak{}rightclick}}-Callback
des anvisierten Nodes auf oder platziert das gehaltene Item, falls es
ein Node ist.

\subsection{on\_drop}\label{on_drop}

\texttt{on\_\allowbreak{}drop}\index{callback@Callback!on_drop@\texttt{on\_\allowbreak{}drop}}
wird aufgerufen, wenn der Spieler ein Item fallen lassen will, etwa mit
der Taste Q oder indem er es aus dem Inventar herauszieht. Standardmäßig
ist das
\texttt{core.\allowbreak{}item\_\allowbreak{}drop}\index{core@\texttt{core}!item_drop@\texttt{core.\allowbreak{}item\_\allowbreak{}drop}},
die das Fallenlassen übernimmt.

\subsection{after\_use}\label{after_use}

\texttt{after\_\allowbreak{}use}\index{callback@Callback!after_use@\texttt{after\_\allowbreak{}use}}
wird beim Abbauen eines Nodes aufgerufen und erlaubt dir, selbst
festzulegen, wie sich ein Werkzeug abnutzt. Fehlt
\texttt{after\_\allowbreak{}use}, entspricht das diesem Code:

\begin{Shaded}
\begin{Highlighting}[]
\VariableTok{after\_use} \OperatorTok{=} \KeywordTok{function}\OperatorTok{(}\VariableTok{itemstack}\OperatorTok{,} \VariableTok{user}\OperatorTok{,} \VariableTok{node}\OperatorTok{,} \VariableTok{digparams}\OperatorTok{)}
    \VariableTok{itemstack}\OperatorTok{:}\NormalTok{add\_wear}\OperatorTok{(}\VariableTok{digparams}\OperatorTok{.}\VariableTok{wear}\OperatorTok{)}
    \ControlFlowTok{return} \VariableTok{itemstack}
\KeywordTok{end}
\end{Highlighting}
\end{Shaded}

\section{item\_place oder
place\_item?}\label{item_place-oder-place_item}

Die Luanti-API bringt viele fertige Callback-Implementierungen mit, die
du verwenden kannst. Sie sind nach dem Schema \emph{Typ zuerst} benannt,
zum Beispiel \texttt{core.\allowbreak{}item\_\allowbreak{}place} und
\texttt{core.\allowbreak{}node\_\allowbreak{}dig}\index{core@\texttt{core}!node_dig@\texttt{core.\allowbreak{}node\_\allowbreak{}dig}}.
Manche werden direkt eingesetzt, andere sind Funktionen, die den
Callback erst zurückgeben:

\begin{Shaded}
\begin{Highlighting}[]
\VariableTok{core}\OperatorTok{.}\NormalTok{register\_item}\OperatorTok{(}\StringTok{"mymod:example"}\OperatorTok{,} \OperatorTok{\{}
    \VariableTok{on\_place} \OperatorTok{=} \VariableTok{core}\OperatorTok{.}\VariableTok{item\_place}\OperatorTok{,}
    \VariableTok{on\_use} \OperatorTok{=} \VariableTok{core}\OperatorTok{.}\NormalTok{item\_eat}\OperatorTok{(}\DecValTok{10}\OperatorTok{),}
\OperatorTok{\})}
\end{Highlighting}
\end{Shaded}

Außerdem enthält die API Funktionen, die etwas \emph{tun}. Sie heißen
verwirrend ähnlich wie die Callback-Implementierungen, beginnen aber mit
dem Verb. Beispiele sind
\texttt{core.\allowbreak{}place\_\allowbreak{}node}\index{core@\texttt{core}!place_node@\texttt{core.\allowbreak{}place\_\allowbreak{}node}}
und
\texttt{core.\allowbreak{}dig\_\allowbreak{}node}\index{core@\texttt{core}!dig_node@\texttt{core.\allowbreak{}dig\_\allowbreak{}node}}:
Mit ihnen kannst du Nodes abbauen und platzieren, mit ähnlicher Wirkung,
als täte es ein Spieler.

\section{Node-Callbacks}\label{node-callbacks}

Solange ein Node in einem Inventar liegt, verwendet er wie oben
beschrieben Item-Callbacks. Ist er in der Welt platziert, verwendet er
Node-Callbacks. Davon gibt es sehr viele, zu viele für dieses Buch;
einige davon kommen aber in späteren Kapiteln vor.

Mehrere dieser Callbacks betreffen Node-Operationen wie das Platzieren
und Entfernen in der Welt. Wichtig: Aus Leistungsgründen werden sie bei
Massenänderungen -- also wenn sehr viele Nodes auf einmal gesetzt werden
-- nicht aufgerufen. Du kannst dich also nicht darauf verlassen, dass
sie immer ausgelöst werden.

\subsection{Rechtsklick und
Platzieren}\label{rechtsklick-und-platzieren}

Klickt der Spieler mit einem Item rechts, während er auf einen Node
zeigt, wird der \texttt{on\_\allowbreak{}place}-Callback des Items
aufgerufen, standardmäßig
\texttt{core.\allowbreak{}item\_\allowbreak{}place}. Hat der anvisierte
Node einen \texttt{on\_\allowbreak{}rightclick}-Callback und hält der
Spieler die Schleichtaste (Shift) \emph{nicht} gedrückt, wird
\texttt{on\_\allowbreak{}rightclick} aufgerufen. Andernfalls platziert
\texttt{core.\allowbreak{}item\_\allowbreak{}place} den Node.

Beim Platzieren eines Nodes werden sowohl
\texttt{on\_\allowbreak{}construct}\index{callback@Callback!on_construct@\texttt{on\_\allowbreak{}construct}}
als auch
\texttt{after\_\allowbreak{}place\_\allowbreak{}node}\index{callback@Callback!after_place_node@\texttt{after\_\allowbreak{}place\_\allowbreak{}node}}
aufgerufen. \texttt{on\_\allowbreak{}construct} wird bei jedem Setzen
eines Nodes außerhalb von Massenänderungen ausgelöst und erhält nur die
Position des Nodes.
\texttt{after\_\allowbreak{}place\_\allowbreak{}node} wird nur beim
Platzieren aufgerufen und bekommt daher mehr Informationen, etwa den
Platzierenden und den Itemstack.\index{itemstack@\texttt{ItemStack}}

Beachte, dass nicht nur Spieler Nodes platzieren können; auch Mobs und
Mods tun das häufig. Deshalb kann \texttt{placer} ein Spieler, eine
Entity\index{entity@Entity} oder \texttt{nil} sein.

\begin{Shaded}
\begin{Highlighting}[]
\VariableTok{core}\OperatorTok{.}\NormalTok{register\_node}\OperatorTok{(}\StringTok{"mymod:mynode"}\OperatorTok{,} \OperatorTok{\{}
    \VariableTok{on\_rightclick} \OperatorTok{=} \KeywordTok{function}\OperatorTok{(}
        \VariableTok{pos}\OperatorTok{,}
        \VariableTok{node}\OperatorTok{,}
        \VariableTok{clicker}\OperatorTok{,}
        \VariableTok{itemstack}\OperatorTok{,}
        \VariableTok{pointed\_thing}
    \OperatorTok{)}
        \ControlFlowTok{if} \VariableTok{clicker}\OperatorTok{:}\NormalTok{is\_player}\OperatorTok{()} \ControlFlowTok{then}
            \VariableTok{core}\OperatorTok{.}\NormalTok{chat\_send\_player}\OperatorTok{(}
                \VariableTok{clicker}\OperatorTok{:}\NormalTok{get\_player\_name}\OperatorTok{(),}
                \StringTok{"Hallo Welt!"}
            \OperatorTok{)}
        \ControlFlowTok{end}
    \KeywordTok{end}\OperatorTok{,}
    \VariableTok{on\_construct} \OperatorTok{=} \KeywordTok{function}\OperatorTok{(}\VariableTok{pos}\OperatorTok{)}
        \KeywordTok{local} \VariableTok{meta} \OperatorTok{=} \VariableTok{core}\OperatorTok{.}\NormalTok{get\_meta}\OperatorTok{(}\VariableTok{pos}\OperatorTok{)}
        \VariableTok{meta}\OperatorTok{:}\NormalTok{set\_string}\OperatorTok{(}\StringTok{"infotext"}\OperatorTok{,} \StringTok{"Mein Node!"}\OperatorTok{)}
    \KeywordTok{end}\OperatorTok{,}
    \VariableTok{after\_place\_node} \OperatorTok{=} \KeywordTok{function}\OperatorTok{(}
        \VariableTok{pos}\OperatorTok{,}
        \VariableTok{placer}\OperatorTok{,}
        \VariableTok{itemstack}\OperatorTok{,}
        \VariableTok{pointed\_thing}
    \OperatorTok{)}
        \CommentTok{{-}{-} placer immer prüfen}
        \ControlFlowTok{if} \VariableTok{placer} \KeywordTok{and} \VariableTok{placer}\OperatorTok{:}\NormalTok{is\_player}\OperatorTok{()} \ControlFlowTok{then}
            \KeywordTok{local} \VariableTok{meta} \OperatorTok{=} \VariableTok{core}\OperatorTok{.}\NormalTok{get\_meta}\OperatorTok{(}\VariableTok{pos}\OperatorTok{)}
            \VariableTok{meta}\OperatorTok{:}\NormalTok{set\_string}\OperatorTok{(}
                \StringTok{"owner"}\OperatorTok{,}
                \VariableTok{placer}\OperatorTok{:}\NormalTok{get\_player\_name}\OperatorTok{()}
            \OperatorTok{)}
        \ControlFlowTok{end}
    \KeywordTok{end}\OperatorTok{,}
\OperatorTok{\})}
\end{Highlighting}
\end{Shaded}

\subsection{Schlagen und Abbauen}\label{schlagen-und-abbauen}

Schlagen heißt, dass der Spieler kurz links klickt. Hat das gehaltene
Item einen \texttt{on\_\allowbreak{}use}-Callback, wird dieser
aufgerufen, andernfalls der \texttt{on\_\allowbreak{}punch}-Callback des
anvisierten Nodes.

Versucht der Spieler, einen Node abzubauen, wird dessen
\texttt{on\_\allowbreak{}dig}\index{callback@Callback!on_dig@\texttt{on\_\allowbreak{}dig}}-Callback
aufgerufen. Standardmäßig ist das
\texttt{core.\allowbreak{}node\_\allowbreak{}dig}: Die Funktion prüft
den Gebietsschutz, nutzt das Werkzeug ab, entfernt den Node und führt
den
\texttt{after\_\allowbreak{}dig\_\allowbreak{}node}\index{callback@Callback!after_dig_node@\texttt{after\_\allowbreak{}dig\_\allowbreak{}node}}-Callback
aus.

\begin{Shaded}
\begin{Highlighting}[]
\VariableTok{core}\OperatorTok{.}\NormalTok{register\_node}\OperatorTok{(}\StringTok{"mymod:mynode"}\OperatorTok{,} \OperatorTok{\{}
    \VariableTok{on\_punch} \OperatorTok{=} \KeywordTok{function}\OperatorTok{(}\VariableTok{pos}\OperatorTok{,} \VariableTok{node}\OperatorTok{,} \VariableTok{puncher}\OperatorTok{,} \VariableTok{pointed\_thing}\OperatorTok{)}
        \ControlFlowTok{if} \VariableTok{puncher}\OperatorTok{:}\NormalTok{is\_player}\OperatorTok{()} \ControlFlowTok{then}
            \VariableTok{core}\OperatorTok{.}\NormalTok{chat\_send\_player}\OperatorTok{(}
                \VariableTok{puncher}\OperatorTok{:}\NormalTok{get\_player\_name}\OperatorTok{(),}
                \StringTok{"Au!"}
            \OperatorTok{)}
        \ControlFlowTok{end}
    \KeywordTok{end}\OperatorTok{,}
\OperatorTok{\})}
\end{Highlighting}
\end{Shaded}

\subsection{\ldots{} und mehr!}\label{und-mehr}

Eine Liste aller Node-Callbacks und weitere Details zu den oben
genannten findest du in der Lua-API-Referenz von Luanti.
